\documentclass[12pt]{article}
\usepackage[utf8]{inputenc}
\usepackage[T1]{fontenc}
\usepackage{hyperref}
\title{Tijaltepec Mixtec Language Resources for K-12 Curriculum Development}
\author{Compiled by Claude Sonnet 5.5 (claude-sonnet-5-5) for Andr\'as Kornai}
\date{}
\begin{document}
\maketitle
\section{Introduction}
Tijaltepec Mixtec (ISO 639-3 xtl, Glottocode tija1235) is an Otomanguean language, also called Mixteco de San Pablo Tijaltepec or Tu'un savi. Wikitongues gives EGIDS 6a (vigorous), while Glottolog lists it as threatened. Little material was found.
\section{Primary Educational Resources}
No primers, textbooks or school programmes were found in the search results.
\section{Linguistic and Reference Materials}
\begin{itemize}
\item Glottolog entry: \url{https://glottolog.org/resource/languoid/id/tija1235}
\item Wikitongues language page: \url{https://wikitongues.org/languages/xtl/}
\end{itemize}
\section{Religious and Community Resources}
\begin{itemize}
\item Scripture Earth listing (notes that speakers may be able to use media in the related San Miguel el Grande or Yosond\'ua Mixtec): \url{https://www.scriptureearth.org/00eng.php?iso=xtl}
\end{itemize}
\section{Digital and Audio Resources}
\begin{itemize}
\item Global Recordings Network, Mixtec, Tijaltepec (audio and slideshow Bible lessons, downloadable MP3 set): \url{https://globalrecordings.net/en/language/xtl}
\end{itemize}
\section{Curriculum Development Context}
The materials found are religious audio and catalogue entries only. Any curriculum would need materials built with the community. Related Mixtec varieties may supply partial models, but intelligibility should not be assumed.
\end{document}